\newpage\begingroup\small\setlength{\emergencystretch}{2em}
\section*{Q4. Generalization and error analysis: final DeBERTa}
\subsection*{Q4.1. Training versus validation}
Training AUROC is 0.995930; fixed-validation AUROC is 0.989423, a gap of 0.006507. The gap is consistent with some overfitting; repeated selection also limits the estimate. Public AUROC falls to 0.92433, showing a larger distribution/generalization gap. These diagnostics use the development checkpoint, not the all-label final fit.
\subsection*{Q4.2. Text-length subgroups}

The same training-derived word-count boundaries (175 and 240 words) define short, medium and long validation groups. AUROC evaluates ranking; the neural classifier's default logit-zero rule corresponds to an uncalibrated softmax value of 0.5.
\par\noindent
\begin{tabularx}{\linewidth}{@{}Yrrr@{}}\toprule
Length group & Count & Neural AUROC & Submitted blend\\\midrule
Long & 646 & 0.991526 & 0.991192 \\
Medium & 643 & 0.987055 & 0.987089 \\
Short & 711 & 0.990983 & 0.986215 \\
\bottomrule\end{tabularx}\par
At the default threshold, the neural model produces 204 false positives and 18 false negatives (88.9\% accuracy). Its high AUROC does not imply a good operating threshold or calibrated probabilities. No validation-optimized threshold is presented as an unbiased test result.

\subsection*{Q4.3. Five false positives and five false negatives}
FP denotes human text predicted AI; FN denotes AI text predicted human. These are five confident errors of each type, not a representative sample. Reasons are inspection-based hypotheses, not proven causes or label corrections.
\par\noindent
\begin{tabularx}{\linewidth}{@{}rrrY@{}}\toprule
ID & Type & Softmax AI & Possible reason from inspected passage\\\midrule
3326 & FP & 0.999685 & Fluent, conventional review structure resembles generated exposition; surface polish is an ambiguous authorship cue.\\
8238 & FP & 0.999608 & Regular academic framing and repeated technical categories may resemble generated abstracts. The first 256 tokens omit later detail.\\
1761 & FP & 0.999542 & Formulaic research-summary language resembles AI academic prose despite concrete method details.\\
3364 & FP & 0.999512 & Structured instructional prose and repeated terms may be mistaken for generated explanations.\\
1585 & FP & 0.999412 & Predictable abstract structure may dominate the decision; input truncation can hide later quantitative details.\\
2209 & FN & 0.008062 & Recipe conventions are shared by both classes; ingredient-heavy opening text supplies weak authorship evidence.\\
7018 & FN & 0.010987 & Technical surface detail resembles a human abstract; authorship classification does not verify the scientific claims.\\
7837 & FN & 0.013223 & Colloquial imperfections and first-person review style resemble human user comments.\\
424 & FN & 0.016915 & Recipe vocabulary and informal wording resemble human recipes, while truncation removes much of the later sequence.\\
7995 & FN & 0.018833 & Reference-style detail and irregular presentation resemble human source text; the model does not verify biographical accuracy.\\
\bottomrule\end{tabularx}\par
\subsection*{Q4.4. Failure patterns and targeted response}
Human abstracts and structured explanations cause confident FP; recipes, informal reviews and biographies cause FN. Shared genres obscure authorship. A controlled 256-versus-384-token comparison with other settings fixed could test whether later context helps. The robustness continuation changes several factors and cannot isolate that hypothesis. Separate domains and calibration need further evaluation.
\endgroup
